\section{The model and the uniform bound}\label{sec:introduction}

For a homogeneous gas, Bose--Einstein condensation means that one orbital
contains a positive fraction of all particles in the equilibrium state.
We consider the spatially constant orbital and keep both the physical
interaction and a positive density fixed as the volume grows.
The bound below uses the same positive fraction throughout a sufficiently
dilute interval of densities.

Let $v:\R^3\to[0,\infty)$ be bounded, measurable, radial, and supported
in a ball of finite radius, with $v$ not zero almost everywhere.
On the torus $\Lambda_L=(\R/L\Z)^3$, of volume $V=L^3$, let
$v_L(x)=\sum_{k\in\Z^3}v(x+Lk)$ and define
\[
 H_{N,L}=-\sum_{i=1}^N\Delta_i+
                \sum_{1\le i<j\le N}v_L(x_i-x_j)
 \quad\hbox{on }L^2_{\mathrm{sym}}(\Lambda_L^N).
\]
The units are $\hbar^2/(2m_{\mathrm{particle}})=k_B=1$.
At $T>0$, the state is the canonical Gibbs state
$\Gamma_{N,L,T}=e^{-H_{N,L}/T}/\Tr e^{-H_{N,L}/T}$; at $T=0$, it is
the ground state. Write $\gamma^{(1)}_{N,L,T}$ for its one-particle
density matrix, normalized to trace $N$, and put $u_0=V^{-1/2}$.

\begin{theorem}\label{thm:main}
For every potential $v$ as above there are constants
$\rho_*(v)>0$ and $c_*(v)>0$ with the following property.
For every $0<\rho<\rho_*(v)$ there is $L_0(\rho,v)<\infty$ such that,
for every $L\ge L_0(\rho,v)$, every integer $N\ge1$ satisfying
$\rho/2\le N/L^3\le2\rho$, and every $0\le T\le\rho^2$,
\[
 \frac{\langle u_0,\gamma^{(1)}_{N,L,T}u_0\rangle}{N}
 \ge c_*(v).
\]
Consequently, for each such fixed density $\rho$ and fixed
$0\le T\le\rho^2$, every sequence with $L\to\infty$ and
$N/L^3\to\rho$ has lower limiting condensate fraction at least $c_*(v)$.
\end{theorem}

The volume threshold may depend on the density; the condensate lower
bound does not. The Hamiltonian is the exact interacting Hamiltonian,
and its potential is not rescaled with the number of particles.
The conclusion is macroscopic occupation of the constant orbital with
a positive lower bound throughout the admitted dilute interval. The
zero-temperature state is unambiguous: positivity of the finite-volume
heat kernel gives a unique bosonic ground state, as shown at the end
of the proof.

\subsection{Context and antecedents}
The ideal-gas theory begins with Bose's counting of light quanta
\cite{Bose1924} and Einstein's extension to a monatomic ideal gas,
including the condensation prediction in his second paper
\cite{Einstein1924IdealGas,Einstein1925IdealGas}.
Bogoliubov's treatment of weakly interacting bosons introduced an
approximate quasiparticle description based on macroscopic occupation of
one mode \cite{Bogolubov1947}.

For interacting particles, Penrose and Onsager formulated condensation
through a macroscopic eigenvalue of the one-particle density matrix and
applied this viewpoint to liquid helium \cite{PenroseOnsager1956}.
Yang developed the broader framework of off-diagonal long-range order,
covering Bose condensation and fermion pairing \cite{Yang1962ODLRO}.
These criteria identify the order to be proved; establishing it for a
specified interacting Hamiltonian remains a separate mathematical problem.

Rigorous lattice antecedents include Dyson, Lieb, and Simon's
low-temperature long-range order for the spin-$\tfrac12$ nearest-neighbor
quantum XY model on the simple cubic lattice in dimensions $d\ge3$
\cite{DysonLiebSimon1978}, and Kennedy, Lieb, and Shastry's ground-state
long-range order for all spins on hypercubic lattices in dimensions $d\ge2$
\cite{KennedyLiebShastry1988XY}.
For spin-$\tfrac12$, the spin--hard-core-boson correspondence relates
these results to condensation of half-filled lattice bosons.
Passing to a continuum Hamiltonian with a fixed interaction
and a prescribed density requires additional arguments.

The distinction between the thermodynamic and Gross--Pitaevskii limits is
essential. Lieb and Seiringer proved complete condensation in the
Gross--Pitaevskii limit, with the interaction scaled so that $Na$ remains
fixed, where $a$ is the scattering length
\cite{LiebSeiringer2002BEC}. Deuchert and Seiringer subsequently treated
positive-temperature homogeneous gases with scattering length
$a_N=a_vL/N$, where $a_v$ is the scattering length of the
unscaled potential, and controlled the one-particle density matrix of approximate
Gibbs minimizers \cite{DeuchertSeiringer2020GP}. These results allow the
interaction scale to vary with particle number. The problem considered
here instead fixes the potential and a positive density before taking
the volume to infinity.

The limits between Gross--Pitaevskii scaling and a fixed-density
thermodynamic limit have several forms. Brennecke, Brooks, Caraci, and
Oldenburg proved complete ground-state condensation in the
Gross--Pitaevskii regime and beyond
\cite{BrenneckeBrooksCaraciOldenburg2025}. Fournais treated low-energy
states on growing boxes in a joint dilute and increasing-volume limit
\cite[Theorem~1.2]{Fournais2021}. Chong, Liang, and Nam gave simpler
proofs of low-energy condensation using kinetic localization and
Poincar\'e-type inequalities \cite{ChongLiangNam2026}.
For the unit-box interaction $N^{2-2\kappa}v(N^{1-\kappa}x)$,
rescaling to a fixed potential gives density $N^{3\kappa-2}$;
the regimes $\kappa<2/3$ therefore couple dilution to increasing
particle number. Junge's low-temperature canonical Gibbs-state
result uses Neumann boxes whose size remains coupled to
the dilute limit \cite[Corollary~6]{Junge2026Propagation}.

The thermodynamic energy problem has a separate history. Dyson's
hard-sphere bounds and the leading dilute asymptotics of Lieb and
Yngvason established the first-order energy scale
\cite{Dyson1957,LiebYngvason1998}. The next correction was predicted
by Lee, Huang, and Yang's pseudopotential calculation
\cite{LeeHuangYang1957}. Fournais and Solovej proved the corresponding
lower bound for nonnegative radial finite-range potentials, including hard cores
\cite[Theorems~1.1 and~1.3]{FournaisSolovej2023}; recent upper bounds
attain the same coefficient for hard spheres and for general
nonnegative even finite-range potentials
\cite{BastiEtAl2026HardSpheres,BastiEtAl2026GeneralPotentials}.
These energy results, and low-temperature free-energy estimates for
strong potentials \cite{FournaisEtAl2026StrongFreeEnergy}, do not by
themselves determine occupation of a single orbital: the one-particle
kinetic gap tends to zero as the volume grows.

A distinct thermodynamic-limit result appears in S\"ut\H{o}'s two-step preprint
\cite{Suto2023TwoStep}. Its Theorem~1.2 assumes nonnegativity both of the
potential and of its Fourier transform, together with integrability,
a Fourier-moment condition, and spatial decay, and states condensation
above a temperature-dependent density threshold. Fourier positivity is
an additional restriction: for example, a nonnegative spherical square
barrier has a Fourier transform that changes sign. That result therefore
does not cover all bounded nonnegative radial finite-range potentials,
and its fixed-temperature threshold does not alone provide the joint
low-density and low-temperature uniformity sought here.

Galanda and Pinamonti construct algebraic equilibrium states with a
condensate background under quantitative convergence conditions for a
spatially localized interaction \cite[Theorem~6.8]{GalandaPinamonti2025Equilibrium}.
Their subsequent analysis treats renormalized correlation functions in
Gross--Pitaevskii scaling, with a zero-temperature limit in its final
convergence theorem \cite[Theorem~7.3]{GalandaPinamonti2026GP}.
These constructions use a different state formulation and coupled
interaction, regulator, and temperature limits from the fixed-potential
canonical thermodynamic limit considered here.

The trajectory formulation descends from Feynman's treatment of the
Bose gas \cite{Feynman1953AtomicLambda}; for the canonical
Feynman--Kac representation and its symmetrization and one-particle
normalizations, see \cite[Appendix~A]{Ueltschi2006FeynmanCycles}.
The random-path comparison has a different ancestry.
Benjamini, Pemantle, and Peres constructed paths with exponentially
decaying intersection tails \cite{BenjaminiPemantlePeres1998}.
H\"aggstr\"om and Mossel developed a hierarchical random-drift
construction with low predictability; its independent block variables
are a close antecedent of the dyadic velocities used below
\cite[Section~3, Proposition~3.1]{HaggstromMossel1998}.
Abbe, Massouli\'e, Montanari, Sly, and Srivastava used finite-endpoint
path averages and overlap estimates for group synchronization
\cite{AbbeMassoulieMontanariSlySrivastava2018}; Garban and Spencer used
this approach for classical disordered-spin symmetry breaking along the Nishimori line
\cite{GarbanSpencer2023}. Here the finite cell paths, their thickened
intersection estimates, and the continuum bridge comparison require
separate proofs.


\subsection{How the comparison works}

Let $q_0=\langle u_0,\gamma^{(1)}u_0\rangle/N$.
The Feynman--Kac formula expresses the partition function by positive
weights on closed collections of paths. Replacing one matching endpoint
by an independently integrated spatial point gives a positive measure
$\mu$ of total mass $q_0$. We construct a measure $\nu$ by changing a
short portion of closed paths, with mass at least a fixed $a>0$, and
show that its density $f=\dd\nu/\dd\mu$ has bounded second moment.
Cauchy--Schwarz then gives
\[
                    a^2\le q_0\int f^2\dd\mu.
\]
Thus the main problem is a comparison of positive path measures.

Partition space into cells of side $s$ comparable to $\rho^{-1/2}$,
and change paths only during a time interval of length $\delta=s^2$.
The mean count $m=Ns^3/V$ is comparable to $\rho^{-1/2}$, making
many particle choices available in a typical cell as the density decreases,
while the slab remains short compared with $1/T$ throughout the
positive-temperature range.

The local analysis must control simultaneous failures in specified
families of cells. Energy and entropy estimates first give occupation
tails and separation of most endpoints at the shorter scale
$r=\zeta\rho^{-1/3}$, with $\zeta>0$ fixed and small. A particle is
eligible when its endpoints are separated, its path stays near them
during short initial and final time intervals, and it satisfies overall
confinement and integrated interaction-score bounds for the path
ensemble. The interacting
bridges are not independent. A separate localized partition-function
comparison bounds their joint trajectory failures, with a cost
proportional to the volume of the tested cells. Together with the
one-time estimates and score bounds, this makes specified sets of
unusable cells unlikely.

Between two distant cells, choose a random lattice path, called a
skeleton, whose encounters with an independent copy have bounded
exponential moments. Here an encounter records a vertex within a
fixed cell distance of the other path; a second estimate weights all
vertices by exponentially decaying distance. These two geometric
bounds are established independently of the gas in
Section~\ref{sec:skeletons}. Detour around unusable components. Select one eligible particle from
each cell on the resulting simple path, and cyclically reassign their
starting points. One end is opened and the selected bridges are sampled
from their exactly normalized interacting law. This defines $\nu$.

The second moment involves two independently chosen paths in related
configurations. Estimating both changes against the original closed
configuration would incur a cost proportional to their lengths.
Instead, a reverse insertion supplies a baseline in which contributions
far from encounters cancel. Two effects remain small as $\rho$ decreases:
an encounter of distinct separated bridges has interaction cost tending
to zero at rate $r^{-1/2}$, and coincident particle choices have probability
of order $1/m$.

Changing trajectories also changes which cells are usable and which
particles are eligible, so the two routes cannot be treated as independent
of the gas. A stricter cell test in the reverse configuration controls
these changes: away from encounters the selection weights agree, while
rare components give a controlled list of possible detours. Keeping
those selection weights throughout the comparison yields the required
second-moment bound after averaging over the environment and the two
paths.

Sections~\ref{sec:energy}--\ref{sec:defects} prove the local analytic
bounds. Section~\ref{sec:geometry} constructs the paths and detours.
Section~\ref{sec:opening} defines the change of measure, and
Section~\ref{sec:comparison} proves its second-moment estimate with all
normalizations and selection constraints retained.
Section~\ref{sec:conclusion} fixes the parameters in order and passes
the uniform positive-temperature bound to zero temperature.
